% --- Table: Semantic layer invariants (auto-fits page width) ---
\begin{table}[htbp]
\centering
\small
\renewcommand{\arraystretch}{1.08}
% squeeze column padding a bit
\setlength{\tabcolsep}{5.5pt}
% final guard: scale whole table if needed
\resizebox{\linewidth}{!}{%
\begin{tabular}{p{2.1cm} p{2cm} p{3cm} p{3cm}}
\toprule
\textbf{Layer} & \textbf{Quantity} & \textbf{Invariant / Monotone} & \textbf{How measured} \\
\midrule
Energy $\mathcal{E}$ &
Lyapunov $E[\psi]$ &
Nonincreasing for neutral windows (A1--A2) &
Energy deltas before/after $\Pi/\Lambda$ pulses; bounded $O(\|A_{\rm obs}\|^2)$ perturbations. \\
\addlinespace[1pt]
Information $\mathcal{I}$ &
Jarlskog $(J_q,J_\ell)$ &
Preserved within tolerances; angles may shift &
$|J_{\rm after}-J_{\rm before}|\le\varepsilon_J$; scaffold logs: quarks $10^{-3}$, leptons $10^{-2}$. \\
\addlinespace[1pt]
Meaning $\mathcal{M}$ &
Identity count $\#\mu$ &
Invariant under legal moves; generation‑lock keeps 3 connected shells &
Counters by fiber filter; minimal $\nu$ filter $\rightarrow$ 24/flavor; “192” superset stable. \\
\bottomrule
\end{tabular}%
}
\caption{Summary of what is preserved, what is monotone, and how it is tracked in code.}
\label{tab:semantic-invariants}
\end{table}
